\documentclass[14pt,a4paper]{extarticle}

\usepackage{fontspec}
\usepackage{polyglossia}
\setdefaultlanguage{russian}
\setotherlanguage{english}
\setmainfont{Times New Roman}
\setmonofont{Courier New}
\newfontfamily\cyrillicfonttt{Times New Roman}

\usepackage[left=30mm,right=15mm,top=20mm,bottom=20mm]{geometry}
\usepackage{setspace}
\onehalfspacing
\setlength{\parindent}{1.25cm}
\setlength{\parskip}{0pt}

\usepackage{xcolor}
\usepackage{amsmath,amssymb}
\usepackage{longtable,booktabs,array,multirow,calc}
\usepackage{graphicx}
\usepackage{etoolbox}
\usepackage{soul}
\usepackage{caption}
\usepackage{hyperref}
\usepackage{xurl}
\usepackage{pdfpages}
\usepackage{listings}
\usepackage{float}

\hypersetup{hidelinks}
\urlstyle{same}
\setcounter{secnumdepth}{3}
\setcounter{tocdepth}{3}
\setlength{\emergencystretch}{3em}
\newcounter{none}

% В pandoc-экспорте встречаются tightlist; оставляем совместимость.
\providecommand{\tightlist}{%
  \setlength{\itemsep}{0pt}\setlength{\parskip}{0pt}}

\captionsetup[figure]{justification=centering,singlelinecheck=false}
\captionsetup[table]{justification=raggedright,singlelinecheck=false}
\captionsetup[figure]{labelsep=endash}
\captionsetup[lstlisting]{position=top,justification=raggedright,singlelinecheck=false,labelsep=endash}
\numberwithin{figure}{section}
\numberwithin{table}{section}
\renewcommand{\lstlistingname}{Листинг}
\makeatletter
\AtBeginDocument{%
  \@ifundefined{c@lstlisting}{}{%
    \@addtoreset{lstlisting}{section}%
    \renewcommand{\thelstlisting}{\thesection.\arabic{lstlisting}}%
  }%
}
\makeatother
\lstset{
  basicstyle=\ttfamily\small,
  frame=single,
  breaklines=true,
  keepspaces=true,
  numbers=left,
  numberstyle=\tiny,
  stepnumber=1,
  numbersep=8pt,
  xleftmargin=1.5em
}

\makeatletter
\patchcmd\longtable{\par}{\if@noskipsec\mbox{}\fi\par}{}{}
\newsavebox\pandoc@box
\newcommand*\pandocbounded[1]{%
  \sbox\pandoc@box{#1}%
  \Gscale@div\@tempa{\textheight}{\dimexpr\ht\pandoc@box+\dp\pandoc@box\relax}%
  \Gscale@div\@tempb{\linewidth}{\wd\pandoc@box}%
  \ifdim\@tempb\p@<\@tempa\p@\let\@tempa\@tempb\fi%
  \ifdim\@tempa\p@<\p@\scalebox{\@tempa}{\usebox\pandoc@box}%
  \else\usebox{\pandoc@box}%
  \fi%
}
\def\fps@figure{htbp}
\makeatother

\begin{document}

\includepdf[pages=-]{docs/Титул.pdf}
\clearpage
\tableofcontents
\clearpage

\section{Общие
положение}\label{ux43eux431ux449ux438ux435-ux43fux43eux43bux43eux436ux435ux43dux438ux435}

\subsection{Анализ предметной области автомобильных
сервисов}\label{ux430ux43dux430ux43bux438ux437-ux43fux440ux435ux434ux43cux435ux442ux43dux43eux439-ux43eux431ux43bux430ux441ux442ux438-ux430ux432ux442ux43eux43cux43eux431ux438ux43bux44cux43dux44bux445-ux441ux435ux440ux432ux438ux441ux43eux432}

Автомобильный сервис (СТО) представляет собой организацию,
осуществляющую техническое обслуживание, диагностику и ремонт
транспортных средств. Взаимодействие с клиентом в автосервисе обычно
начинается на этапе первичного обращения, когда владелец автомобиля
сообщает о проблеме, симптомах неисправности или планируемых работах. На
данном этапе происходит сбор исходной информации, необходимой для
последующей диагностики и планирования обслуживания, а также формируется
предварительное представление о характере обращения.

Первичная консультация клиента включает уточнение сведений об автомобиле
(марка, модель, год выпуска, пробег), условий эксплуатации,
обстоятельств появления неисправности, а также описание симптомов в
свободной форме (например, «стук в подвеске на кочках», «плавают
обороты», «горит чек»). Качество и полнота собранной информации напрямую
влияет на точность предварительных рекомендаций и дальнейшую работу
менеджера сервиса.

В современных автосервисах приём обращений клиентов осуществляется через
различные каналы: телефонные звонки, мессенджеры, формы на сайте,
социальные сети, а также личное обращение. Указанные каналы формируют
поток обращений различного качества: часть клиентов предоставляет точные
и структурированные данные, однако значительная доля обращений содержит
неполные, противоречивые или эмоционально окрашенные сообщения, что
усложняет работу менеджера и увеличивает время обработки заявки.

В рамках предметной области выделяются основные участники процесса
первичной консультации: клиент автосервиса, менеджер (администратор
приёма), мастер-приёмщик, а также технические специалисты (диагносты,
механики). Менеджер выполняет роль первичного фильтра: принимает
обращение, задаёт уточняющие вопросы, фиксирует данные в информационной
системе или вручную, формирует заявку и передаёт её в производственный
контур автосервиса.

При отсутствии автоматизации процесс первичной консультации
характеризуется высокой зависимостью от человеческого фактора. Разные
сотрудники могут по-разному формулировать уточняющие вопросы,
фиксировать данные с различной полнотой и допускать ошибки при
интерпретации обращений. Дополнительной сложностью является то, что
клиенты часто используют бытовые выражения и неточные термины для
описания неисправностей, что требует от менеджера опыта и знаний для
правильной интерпретации.

Таким образом, предметная область первичной консультации автосервиса
включает процессы приёма обращения, уточнения входной информации,
предварительного определения вида услуги и формирования заявки.
Автоматизация данного этапа позволяет унифицировать сбор данных,
сократить время первичной обработки обращений и повысить качество
информации, передаваемой специалистам для дальнейшего обслуживания.

\subsection{Постановка проблемы и выявление недостатков
существующего
процесса}\label{ux43fux43eux441ux442ux430ux43dux43eux432ux43aux430-ux43fux440ux43eux431ux43bux435ux43cux44b-ux438-ux432ux44bux44fux432ux43bux435ux43dux438ux435-ux43dux435ux434ux43eux441ux442ux430ux442ux43aux43eux432-ux441ux443ux449ux435ux441ux442ux432ux443ux44eux449ux435ux433ux43e-ux43fux440ux43eux446ux435ux441ux441ux430}

Несмотря на широкое распространение информационных систем в сфере
автомобильного сервиса, этап первичной консультации клиентов во многих
организациях по-прежнему реализуется преимущественно вручную. Обработка
обращений осуществляется менеджерами или администраторами, которые
принимают сообщения по телефону, через мессенджеры или формы обратной
связи и самостоятельно уточняют детали неисправности.

Анализ существующего процесса показал, что ручная обработка первичных
обращений сопровождается рядом проблем.

Во-первых, наблюдается высокая нагрузка на сотрудников автосервиса. В
часы пик количество входящих обращений значительно возрастает, что
приводит к увеличению времени ожидания ответа и снижению качества
консультации. Менеджеры вынуждены параллельно обрабатывать несколько
запросов, что увеличивает вероятность ошибок при фиксации данных.

Во-вторых, отсутствует единый стандарт сбора информации. Разные
сотрудники задают разные уточняющие вопросы, что приводит к
неоднородности данных. В одних случаях информация фиксируется подробно
(марка автомобиля, пробег, условия проявления неисправности), в других
-- ограничивается общим описанием проблемы. Это затрудняет дальнейшую
диагностику и планирование работ.

В-третьих, качество первичной консультации напрямую зависит от опыта
конкретного менеджера. Недостаточная квалификация или усталость
сотрудника могут привести к некорректной интерпретации обращения
клиента. Кроме того, клиенты часто используют бытовые формулировки,
которые требуют профессионального анализа для корректного понимания.

В-четвёртых, процесс первичной консультации не всегда обеспечивает
структурированное хранение данных. Информация может сохраняться в разных
системах (CRM, мессенджеры, таблицы), что усложняет последующий анализ
обращений и формирование статистики.

Таким образом, выявленные недостатки существующего процесса первичной
консультации можно систематизировать следующим образом:


В рамках рассматриваемого вопроса выделяются следующие положения: во-первых, высокая нагрузка на персонал при увеличении количества обращений; во-вторых, зависимость качества консультации от человеческого фактора; в-третьих, отсутствие унифицированного алгоритма сбора информации; в-четвертых, неструктурированность входных данных; в-пятых, увеличение времени обработки обращения клиента.


Указанные проблемы обосновывают необходимость разработки
интеллектуальной системы, обеспечивающей автоматизированный сбор
информации, анализ текстовых обращений и формирование структурированных
заявок на обслуживание. Внедрение ИИ-ассистента на этапе первичной
консультации позволяет стандартизировать диалог, повысить полноту и
точность собираемых данных, а также снизить нагрузку на сотрудников
автомобильного сервиса.

\subsection{Описание целевой аудитории
системы}\label{ux43eux43fux438ux441ux430ux43dux438ux435-ux446ux435ux43bux435ux432ux43eux439-ux430ux443ux434ux438ux442ux43eux440ux438ux438-ux441ux438ux441ux442ux435ux43cux44b}

Разрабатываемая интеллектуальная система ориентирована на несколько
категорий пользователей, каждая из которых имеет собственные цели и
сценарии взаимодействия с системой.

Клиенты автомобильного сервиса

Основной целевой аудиторией являются владельцы транспортных средств,
обращающиеся в автосервис по вопросам диагностики, ремонта или
технического обслуживания автомобиля.

Для данной категории пользователей характерны следующие особенности:


В рамках рассматриваемого вопроса выделяются следующие положения: во-первых, отсутствие специализированных технических знаний; во-вторых, использование бытовых формулировок при описании неисправностей; в-третьих, необходимость быстрого получения предварительной информации о характере проблемы и ориентировочной стоимости услуг; в-четвертых, потребность в удобном и доступном формате взаимодействия (веб-интерфейс, мобильный доступ).


Клиент заинтересован в оперативной консультации без необходимости
ожидания ответа менеджера, а также в прозрачности процесса формирования
заявки.

Менеджеры автомобильного сервиса

Менеджеры (администраторы приёма) являются внутренними пользователями
системы. Их задача -- обработка заявок, взаимодействие с клиентами и
координация дальнейших действий.

Для данной категории характерны:


В рамках рассматриваемого вопроса выделяются следующие положения: во-первых, Высокая нагрузка при большом количестве обращений; во-вторых, необходимость оперативной обработки заявок; в-третьих, потребность в структурированной и полной информации о проблеме автомобиля; в-четвертых, заинтересованность в снижении времени первичной консультации.


Использование интеллектуальной системы позволяет менеджеру получать уже
структурированную заявку с предварительными выводами, что упрощает
дальнейшую работу.

Руководство автомобильного сервиса

Руководители заинтересованы в:


В рамках рассматриваемого вопроса выделяются следующие положения: во-первых, повышении эффективности работы сотрудников; во-вторых, сокращении времени обработки обращений; в-третьих, получении аналитических данных о типах неисправностей и загруженности сервиса; в-четвертых, повышении уровня удовлетворённости клиентов.


Для данной категории важна возможность анализа статистики обращений и
контроля эффективности работы системы.

Администратор системы

Системный администратор отвечает за настройку, сопровождение и
мониторинг работы системы.

Его задачи включают:


В рамках рассматриваемого вопроса выделяются следующие положения: во-первых, управление ролями пользователей; во-вторых, настройку сценариев консультации; в-третьих, контроль корректности работы ИИ-модуля; в-четвертых, обеспечение информационной безопасности и устойчивости системы.


Таким образом, разрабатываемая система ориентирована как на внешних
пользователей (клиентов), так и на внутренних сотрудников автосервиса.
Каждая категория пользователей взаимодействует с системой в рамках
собственной роли, что требует реализации ролевой модели доступа и
разграничения функциональных возможностей.

\subsection{Моделирование бизнес-процессов первичной консультации
клиентов}\label{ux43cux43eux434ux435ux43bux438ux440ux43eux432ux430ux43dux438ux435-ux431ux438ux437ux43dux435ux441-ux43fux440ux43eux446ux435ux441ux441ux43eux432-ux43fux435ux440ux432ux438ux447ux43dux43eux439-ux43aux43eux43dux441ux443ux43bux44cux442ux430ux446ux438ux438-ux43aux43bux438ux435ux43dux442ux43eux432}

Для обоснования необходимости разработки интеллектуальной системы было
выполнено моделирование бизнес-процесса первичной консультации клиентов
автомобильного сервиса в двух состояниях: текущем (AS-IS) и целевом
(TO-BE). Такой подход позволяет наглядно продемонстрировать недостатки
существующей модели и определить направления её совершенствования.

Модель текущего процесса (AS-IS)

В текущем состоянии процесс первичной консультации реализуется
преимущественно вручную. Клиент обращается в автосервис через доступные
каналы коммуникации -- телефон, мессенджеры или веб-форму. Менеджер
принимает обращение, уточняет сведения об автомобиле и характере
неисправности, после чего фиксирует полученную информацию в CRM-системе
либо вручную. При недостаточности данных менеджер повторно связывается с
клиентом для уточнения деталей. На основании собранной информации
формируется заявка на обслуживание и передаётся специалистам.

\begin{figure}[H]
\centering
\includegraphics[width=0.95\textwidth]{thesis/images/media/image1.jpg}
\caption{BPMN-диаграмма текущего процесса первичной консультации (AS-IS)}
\label{fig:asis-bpmn}
\end{figure}

Данный процесс характеризуется высокой зависимостью от компетенции
сотрудника и не имеет строго регламентированного алгоритма уточнения
данных. Описание неисправности часто формулируется клиентом в свободной
форме, что требует дополнительной интерпретации. В условиях высокой
нагрузки увеличивается вероятность неполного сбора информации, что может
привести к задержкам на последующих этапах обслуживания. Кроме того,
предварительная оценка стоимости работ носит субъективный характер и
основывается на опыте конкретного менеджера.

Модель целевого процесса (TO-BE)

Целевая модель процесса предполагает внедрение интеллектуальной системы
первичной консультации. В этом случае взаимодействие начинается с
диалогового интерфейса, через который клиент описывает проблему
автомобиля. Система анализирует текстовое сообщение, выделяет ключевые
параметры обращения и при необходимости формирует уточняющие вопросы.
После получения достаточного объёма данных выполняется классификация
обращения по видам услуг, формируются предварительные рекомендации и
ориентировочная стоимость обслуживания в формате «от». Дополнительно
рассчитывается показатель уверенности в корректности сформированного
результата. На основе полученных данных автоматически создаётся
структурированная заявка, которая передаётся менеджеру автосервиса для
дальнейшей работы.

Для детализации взаимодействия компонентов системы построена
UML-диаграмма последовательности.

\begin{figure}[H]
\centering
\includegraphics[width=0.9\textwidth]{thesis/images/media/image2.png}
\caption{UML-диаграмма последовательности процесса первичной консультации и формирования заявки}
\label{fig:sequence-uml}
\end{figure}

Диаграмма последовательности отражает порядок взаимодействия клиента,
веб-интерфейса, серверной части, локального модуля искусственного
интеллекта, базы данных и канала уведомлений в процессе первичной
консультации и формирования заявки.

Процесс начинается с инициации консультации клиентом через
веб-интерфейс. Frontend-компонент отправляет запрос на создание
консультационной сессии в серверную часть, после чего информация
сохраняется в базе данных. Далее пользователь вводит описание проблемы,
которое передаётся на сервер и сохраняется в контексте текущей сессии.

Серверная часть передаёт текст обращения в модуль искусственного
интеллекта для анализа. Интеллектуальный модуль формирует
структурированный результат, включающий возможные гипотезы
неисправностей, список недостающих данных и показатель уверенности
(confidence).

В альтернативном фрагменте (alt) показаны два возможных сценария
развития процесса:


В рамках рассматриваемого вопроса выделяются следующие положения: во-первых, при недостаточности данных система формирует уточняющий вопрос и возвращает его пользователю; во-вторых, при достаточности данных выполняется сохранение результатов анализа, расчёт ориентировочной стоимости услуг в формате «от» и отображение предварительного результата клиенту.


После подтверждения пользователем создания заявки сервер формирует
запись в базе данных и инициирует передачу уведомления менеджеру через
интеграцию с Telegram-ботом. Завершается процесс уведомлением клиента об
успешном создании заявки.

Структурная модель данных

Для формализации логики хранения информации и обеспечения корректной
работы системы была разработана структурная модель данных, отражающая
взаимосвязь пользователей, консультационных сессий, сообщений и
сформированных заявок.

\begin{figure}[H]
\centering
\includegraphics[height=0.7\textheight]{thesis/images/media/image3.jpeg}
\caption{Схема работы консультации интеллектуальной системы}
\label{fig:consult-logic}
\end{figure}

Структурная модель демонстрирует, что каждая консультационная сессия
связана с пользователем, содержит историю сообщений и может завершаться
созданием заявки на обслуживание. Такое представление обеспечивает
целостность данных, возможность хранения истории взаимодействия и
последующего анализа обращений.

Таким образом, моделирование бизнес-процессов и структуры данных
позволяет обосновать переход от ручной схемы обработки обращений к
интеллектуализированной системе, обеспечивающей автоматизированный
анализ, формирование предварительных выводов и создание
структурированной заявки.

\subsection{Анализ аналогов и существующих
решений}\label{ux430ux43dux430ux43bux438ux437-ux430ux43dux430ux43bux43eux433ux43eux432-ux438-ux441ux443ux449ux435ux441ux442ux432ux443ux44eux449ux438ux445-ux440ux435ux448ux435ux43dux438ux439}

В целях обоснования необходимости разработки интеллектуальной системы
первичной консультации клиентов автомобильного сервиса был проведён
анализ существующих решений, применяемых в сфере автоматизации
взаимодействия с клиентами.

В настоящее время в автомобильном сервисе используются следующие
категории программных продуктов:


В рамках рассматриваемого вопроса выделяются следующие положения: во-первых, CRM-системы для управления заявками и клиентской базой.; во-вторых, Онлайн-формы обратной связи на сайтах автосервисов.; в-третьих, Чат-боты на основе заранее заданных сценариев.; в-четвертых, Мессенджеры и социальные сети как каналы приёма обращений.


CRM-системы (Bitrix24, AmoCRM и др.) обеспечивают хранение информации о
клиентах и заявках, однако не выполняют интеллектуального анализа
текстового обращения. Данные в систему вносятся сотрудником вручную, что
не устраняет проблему человеческого фактора и не автоматизирует этап
уточнения информации.

Формы обратной связи на сайтах позволяют структурировать часть входных
данных (например, указание марки автомобиля и контактных данных), однако
не обеспечивают адаптивного диалога. Пользователь самостоятельно
заполняет поля формы, и при недостаточности информации менеджеру
приходится дополнительно связываться с клиентом для уточнения деталей.

Чат-боты, реализованные на основе заранее запрограммированных сценариев,
позволяют частично автоматизировать процесс консультации. Однако такие
решения работают по жёстко заданным ветвям диалога и не обладают
возможностью полноценного анализа свободного текстового описания
проблемы. При выходе пользователя за пределы предусмотренного сценария
система не способна корректно интерпретировать обращение.

Использование мессенджеров и социальных сетей не предполагает
автоматизированной обработки данных и полностью зависит от действий
сотрудника автосервиса.

Таким образом, существующие решения либо выполняют функцию хранения и
маршрутизации заявок, либо обеспечивают простую сценарную автоматизацию
без применения методов обработки естественного языка. Они не реализуют
комплексный интеллектуальный анализ текстового описания неисправности
автомобиля и не формируют вероятностную оценку возможных причин
обращения.

Разрабатываемая система отличается от рассмотренных аналогов тем, что
включает подсистему интеллектуальной обработки обращений, способную:


В рамках рассматриваемого вопроса выделяются следующие положения: во-первых, анализировать текстовые сообщения клиента в свободной форме; во-вторых, извлекать структурированные параметры обращения (марка и модель автомобиля, год выпуска, пробег, симптомы, условия проявления неисправности); в-третьих, формировать уточняющие вопросы при недостаточности или противоречивости данных; в-четвертых, выполнять классификацию обращения по видам услуг автомобильного сервиса; в-пятых, рассчитывать показатель уверенности (confidence) в корректности сформированных предварительных выводов; в-шестых, формировать ориентировочную предварительную стоимость услуг в формате «от», на основе типа предполагаемой неисправности и выбранного вида обслуживания; в-седьмых, автоматически создавать структурированную заявку на обслуживание для передачи менеджеру автосервиса.


Таким образом, разрабатываемая интеллектуальная система обеспечивает не
только автоматизацию сбора и хранения информации, но и формирование
предварительной аналитической оценки обращения, включая вероятностную
диагностику и ориентировочную стоимость услуг, что позволяет повысить
прозрачность процесса обслуживания и сократить время взаимодействия
между клиентом и автомобильным сервисом.

\subsection{Формирование требований к интеллектуальной
системе}\label{ux444ux43eux440ux43cux438ux440ux43eux432ux430ux43dux438ux435-ux442ux440ux435ux431ux43eux432ux430ux43dux438ux439-ux43a-ux438ux43dux442ux435ux43bux43bux435ux43aux442ux443ux430ux43bux44cux43dux43eux439-ux441ux438ux441ux442ux435ux43cux435}

На основании проведённого анализа предметной области, изучения
существующих решений и моделирования бизнес-процессов была сформирована
совокупность требований к разрабатываемой интеллектуальной системе
первичной консультации клиентов автомобильного сервиса.

Разрабатываемая система должна обеспечивать диалоговое взаимодействие с
пользователем в режиме реального времени, поддерживать последовательный
сбор информации о транспортном средстве и симптомах неисправности, а
также выполнять интеллектуальный анализ текстовых обращений. Система
должна быть способна выявлять недостающие сведения и формировать
уточняющие вопросы до момента получения достаточного объёма данных для
предварительного вывода.

К функциональным требованиям системы относятся:

анализ текстового обращения пользователя с использованием модели
обработки естественного языка;

извлечение структурированных параметров (марка автомобиля, модель, год
выпуска, симптомы, условия проявления неисправности);


В рамках рассматриваемого вопроса выделяются следующие положения: во-первых, формирование вероятностного перечня возможных причин неисправности; во-вторых, расчёт показателя уверенности (confidence) на основе полноты и согласованности данных; в-третьих, определение рекомендуемого вида услуги; в-четвертых, формирование ориентировочной стоимости обслуживания в формате «от» на основании справочника услуг; в-пятых, автоматическое создание структурированной заявки и передача её менеджеру автосервиса; в-шестых, хранение истории консультации и результатов анализа в базе данных.


К нефункциональным требованиям относятся:


В рамках рассматриваемого вопроса выделяются следующие положения: во-первых, автономность работы интеллектуального модуля (локальное развёртывание модели без зависимости от сторонних платных сервисов); во-вторых, обеспечение конфиденциальности пользовательских данных; в-третьих, корректное функционирование на двух платформах -- персональных компьютерах и мобильных устройствах; в-четвертых, адаптивность пользовательского интерфейса; в-пятых, масштабируемость серверной части при увеличении числа обращений; в-шестых, устойчивость к ошибочным или некорректным входным данным.


Особое значение имеет требование к прозрачности работы интеллектуального
модуля. Пользователь должен получать не только итоговую рекомендацию, но
и понимать предварительный характер результата. По этой причине система
обязана отображать вероятностную оценку гипотез и ориентировочную
стоимость услуг как предварительную информацию, не заменяющую
полноценную техническую диагностику.

Таким образом, сформулированные требования отражают как функциональную
направленность системы на интеллектуальную обработку обращений, так и
организационные и технические условия её эксплуатации в среде
автомобильного сервиса.

\subsection{Обзор современных технологий и платформ для реализации
интеллектуальной
системы}\label{ux43eux431ux437ux43eux440-ux441ux43eux432ux440ux435ux43cux435ux43dux43dux44bux445-ux442ux435ux445ux43dux43eux43bux43eux433ux438ux439-ux438-ux43fux43bux430ux442ux444ux43eux440ux43c-ux434ux43bux44f-ux440ux435ux430ux43bux438ux437ux430ux446ux438ux438-ux438ux43dux442ux435ux43bux43bux435ux43aux442ux443ux430ux43bux44cux43dux43eux439-ux441ux438ux441ux442ux435ux43cux44b}

Разработка интеллектуальной системы первичной консультации клиентов
автомобильного сервиса предполагает использование современных
технологий, обеспечивающих диалоговую обработку текстовых обращений,
автономность работы и возможность функционирования на нескольких
платформах. В рамках проекта ключевым элементом является подсистема
искусственного интеллекта, а веб-интерфейс рассматривается как средство
взаимодействия пользователя с интеллектуальным модулем.

Современные интеллектуальные системы диалогового типа реализуются на
основе моделей обработки естественного языка (Natural Language
Processing, NLP) и больших языковых моделей (LLM -- Large Language
Models). Такие модели способны анализировать текстовые сообщения,
извлекать сущности (марка автомобиля, симптомы, условия проявления
неисправности), выявлять логические связи между признаками и формировать
уточняющие вопросы в процессе диалога. В отличие от статических
чат-ботов с жёстко заданными сценариями, LLM-модели обеспечивают
адаптивное взаимодействие и позволяют учитывать контекст предыдущих
сообщений пользователя.

С учётом требований автономности и отсутствия зависимости от внешних
платных сервисов в рамках разрабатываемой системы предполагается
локальное развёртывание языковой модели на серверной инфраструктуре.
Такой подход позволяет:


В рамках рассматриваемого вопроса выделяются следующие положения: во-первых, исключить использование сторонних API-ключей; во-вторых, обеспечить контроль над обработкой данных; в-третьих, повысить уровень конфиденциальности пользовательской информации; в-четвертых, обеспечить стабильность работы системы вне зависимости от внешних сервисов.


Интеграция интеллектуального модуля осуществляется через серверную часть
системы. Веб-интерфейс передаёт текст сообщения пользователя в серверное
приложение, которое сохраняет данные в рамках консультационной сессии и
передаёт их в ИИ-модуль. Модель анализирует сообщение, формирует
структурированный результат (извлечённые параметры, гипотезы, уточняющие
вопросы, показатель уверенности), после чего сервер обрабатывает
полученные данные и возвращает ответ пользователю.

Архитектурно система реализуется по клиент-серверной модели. Клиентская
часть представляет собой адаптивное веб-приложение, обеспечивающее
корректное отображение интерфейса как на персональных компьютерах, так и
на мобильных устройствах. Для обеспечения работы на двух платформах
используется подход Progressive Web Application (PWA), позволяющий
запускать систему в браузере на ПК и устанавливать её как веб-приложение
на смартфоне без разработки отдельного нативного мобильного клиента.

Серверная часть реализуется на платформе Node.js, обеспечивающей
асинхронную обработку запросов и удобную интеграцию с модулями машинного
обучения. Для взаимодействия с базой данных используется PostgreSQL,
позволяющая хранить структурированные данные о пользователях,
консультационных сессиях, истории диалога, результатах анализа и
сформированных заявках.

Формирование предварительной стоимости услуг в формате «от» реализуется
на основе связки интеллектуального анализа и справочника услуг.
ИИ-модуль определяет наиболее вероятный тип неисправности и
рекомендуемый вид обслуживания, после чего серверная часть сопоставляет
результат с таблицей услуг и диапазоном цен, хранящихся в базе данных.
Таким образом, расчёт стоимости осуществляется контролируемым алгоритмом
и не зависит напрямую от модели.

Показатель уверенности (confidence) формируется на основе полноты
собранных данных, согласованности симптомов и вероятностной оценки
гипотез модели. Это позволяет количественно оценить достоверность
предварительного вывода и повысить прозрачность работы системы.

Таким образом, выбранный технологический стек обеспечивает автономность
интеллектуального модуля, возможность диалогового взаимодействия с
пользователем, формирование вероятностных выводов и предварительной
стоимости услуг, а также функционирование системы на двух платформах --
персональных компьютерах и мобильных устройствах -- без разработки
отдельного мобильного приложения.

\subsection{Проектирование пользовательского интерфейса
системы}\label{ux43fux440ux43eux435ux43aux442ux438ux440ux43eux432ux430ux43dux438ux435-ux43fux43eux43bux44cux437ux43eux432ux430ux442ux435ux43bux44cux441ux43aux43eux433ux43e-ux438ux43dux442ux435ux440ux444ux435ux439ux441ux430-ux441ux438ux441ux442ux435ux43cux44b}

Проектирование пользовательского интерфейса интеллектуальной системы
первичной консультации направлено на обеспечение интуитивно понятного
взаимодействия клиента с ИИ-ассистентом и функциональными возможностями
сервиса. Интерфейс реализован в виде адаптивного веб-приложения,
обеспечивающего корректную работу как на персональных компьютерах, так и
на мобильных устройствах.

Для входа в систему или регистрации пользователю необходимо
воспользоваться соответствующими кнопками, расположенными в верхней
части веб-приложения.

\begin{figure}[H]
\centering
\includegraphics[width=0.9\textwidth]{thesis/images/media/image4.png}
\caption{Кнопки перехода к регистрации и авторизации}
\label{fig:auth-buttons}
\end{figure}

При нажатии на кнопку «Регистрация» открывается модальное окно с формой
создания учетной записи клиента. Пользователь указывает адрес
электронной почты и пароль, которые в дальнейшем используются для
авторизации в системе. Закрытие окна осуществляется нажатием кнопки
закрытия, расположенной в правом верхнем углу модального окна.

\begin{figure}[H]
\centering
\includegraphics[width=0.9\textwidth]{thesis/images/media/image5.png}
\caption{Процесс регистрации пользователя}
\label{fig:register-process}
\end{figure}

При выборе функции «Войти» открывается форма авторизации, в которой
пользователь вводит зарегистрированные учетные данные. После успешной
проверки введённой информации система предоставляет доступ к функционалу
первичной консультации.

\begin{figure}[H]
\centering
\includegraphics[width=0.45\textwidth]{thesis/images/media/image6.png}
\caption{Процесс авторизации пользователя}
\label{fig:login-process}
\end{figure}

После авторизации пользователь перенаправляется в основной рабочий
интерфейс системы -- страницу диалоговой консультации. Интерфейс
реализован в формате чата, где отображаются сообщения пользователя и
ИИ-ассистента, статус текущей консультационной сессии, а также поле
ввода текстовых сообщений.

\begin{figure}[H]
\centering
\includegraphics[width=0.85\textwidth]{thesis/images/media/image7.png}
\caption{Страница диалогового интерфейса системы}
\label{fig:dialog-page}
\end{figure}

В процессе консультации система формирует уточняющие вопросы на основе
ранее полученной информации. Диалог строится последовательно, что
позволяет структурировать входные данные и повышать точность
предварительных выводов.

\begin{figure}[H]
\centering
\includegraphics[width=0.75\textwidth]{thesis/images/media/image8.png}
\caption{Страница уточняющих вопросов ассистента}
\label{fig:clarify-page}
\end{figure}

После получения достаточного объёма информации система формирует
предварительные результаты анализа обращения. В соответствующем разделе
интерфейса отображаются предполагаемый вид услуги, ориентировочная
стоимость обслуживания в формате «от», вероятностная оценка возможных
неисправностей и рекомендации по дальнейшим действиям. Пользователь
может подтвердить создание заявки или инициировать обратный звонок
менеджера.

\begin{figure}[H]
\centering
\includegraphics[width=0.95\textwidth]{thesis/images/media/image9.png}
\caption{Вкладка предварительных результатов консультации}
\label{fig:preliminary-results}
\end{figure}

В личном кабинете клиент получает доступ к истории обращений и
сформированных заявок. В данном разделе отображаются темы консультаций,
даты обращений и текущие статусы обработки заявок. Это обеспечивает
прозрачность взаимодействия с сервисом и возможность отслеживания этапов
обслуживания.

\begin{figure}[H]
\centering
\includegraphics[width=0.9\textwidth]{thesis/images/media/image9.png}
\caption{Страница личного кабинета пользователя}
\label{fig:user-account}
\end{figure}

Таким образом, интерфейс системы обеспечивает последовательный переход
от регистрации пользователя к диалоговой консультации, получению
предварительного результата и формированию заявки. Структура интерфейса
ориентирована на упрощение взаимодействия клиента с интеллектуальным
ассистентом и минимизацию количества действий, необходимых для подачи
обращения в автомобильный сервис.

\subsection{Выводы по
главе}\label{ux432ux44bux432ux43eux434ux44b-ux43fux43e-ux433ux43bux430ux432ux435}

В первой главе выпускной квалификационной работы был проведён
комплексный анализ предметной области автомобильного сервиса в части
организации первичной консультации клиентов. Выявлено, что существующий
процесс обработки обращений характеризуется высокой зависимостью от
человеческого фактора, отсутствием формализованного алгоритма уточнения
данных и субъективностью предварительной оценки стоимости работ.

Проведён анализ существующих программных решений, используемых в
автосервисах, включая CRM-системы и сценарные чат-боты. Установлено, что
данные инструменты не обеспечивают полноценной интеллектуальной
обработки текстовых обращений и не формируют вероятностную оценку
неисправностей на этапе первичного взаимодействия с клиентом.

В результате моделирования бизнес-процессов в состояниях AS-IS и TO-BE
обоснована целесообразность внедрения интеллектуального ассистента,
выполняющего автоматизированный анализ обращений, формирование
уточняющих вопросов, расчёт показателя уверенности и ориентировочной
стоимости услуг в формате «от». Разработана контекстная и процессная
модель системы, отражающая взаимодействие пользователей, серверной части
и модуля искусственного интеллекта.

Сформированы функциональные и нефункциональные требования к
разрабатываемой системе, включая требования к автономности работы
ИИ-модуля, адаптивности интерфейса и обеспечению безопасности данных.

Таким образом, результаты проведённого анализа подтверждают актуальность
разработки интеллектуальной системы первичной консультации клиентов
автосервиса и создают основу для перехода к этапу проектирования
архитектуры системы и описанию алгоритмов её функционирования, что будет
рассмотрено во второй главе.

\section{\texorpdfstring{\hfill\break
Техническая
реализация}{ Техническая реализация}}\label{ux442ux435ux445ux43dux438ux447ux435ux441ux43aux430ux44f-ux440ux435ux430ux43bux438ux437ux430ux446ux438ux44f}

\subsection{Проектирование архитектуры автоматизированной системы
первичной консультации клиентов автомобильного
сервиса}\label{ux43fux440ux43eux435ux43aux442ux438ux440ux43eux432ux430ux43dux438ux435-ux430ux440ux445ux438ux442ux435ux43aux442ux443ux440ux44b-ux430ux432ux442ux43eux43cux430ux442ux438ux437ux438ux440ux43eux432ux430ux43dux43dux43eux439-ux441ux438ux441ux442ux435ux43cux44b-ux43fux435ux440ux432ux438ux447ux43dux43eux439-ux43aux43eux43dux441ux443ux43bux44cux442ux430ux446ux438ux438-ux43aux43bux438ux435ux43dux442ux43eux432-ux430ux432ux442ux43eux43cux43eux431ux438ux43bux44cux43dux43eux433ux43e-ux441ux435ux440ux432ux438ux441ux430}

Во второй главе выпускной квалификационной работы рассматриваются
вопросы технической реализации автоматизированной системы первичной
консультации клиентов автомобильного сервиса на основе технологий
искусственного интеллекта.

На основании результатов анализа предметной области, выполненного в
первой главе, сформированы требования к архитектуре программного
решения, структуре базы данных, пользовательскому интерфейсу и
подсистеме интеллектуальной обработки текстовых обращений пользователей.
В соответствии с поставленными задачами разработана клиент-серверная
архитектура системы, обеспечивающая взаимодействие веб-интерфейса
пользователя, серверной части приложения, базы данных и локально
развёрнутой языковой модели.

В рамках главы выполнено проектирование структуры программного
обеспечения, определён технологический стек разработки, разработана
модель данных приложения, реализованы клиентская и серверная части
системы, а также обеспечена интеграция подсистемы искусственного
интеллекта, предназначенной для анализа пользовательских обращений и
формирования предварительных диагностических рекомендаций.

Особое внимание уделено реализации диалогового интерфейса консультации,
обеспечивающего последовательный сбор диагностически значимых параметров
обращения пользователя, а также разработке алгоритма формирования
структурированного диагностического заключения на основе гибридного
подхода, объединяющего rule-based методы и возможности языковой модели.

Таким образом, во второй главе рассматриваются основные этапы
проектирования и программной реализации автоматизированной системы
первичной консультации клиентов автомобильного сервиса, обеспечивающей
повышение эффективности обработки обращений пользователей и сокращение
времени первичной диагностики неисправностей транспортных средств.

\subsection{Архитектура
приложения}\label{ux430ux440ux445ux438ux442ux435ux43aux442ux443ux440ux430-ux43fux440ux438ux43bux43eux436ux435ux43dux438ux44f}

Разрабатываемая автоматизированная информационная система первичной
консультации клиентов автомобильного сервиса реализована в виде
веб-приложения клиент-серверного типа с интегрированной подсистемой
интеллектуальной обработки пользовательских обращений. Архитектура
системы построена с учётом требований автономности функционирования
языковой модели, масштабируемости программной структуры и обеспечения
последовательного диалогового взаимодействия пользователя с системой.

В рамках разработанного решения применяется модульный подход к
построению архитектуры приложения, предусматривающий разделение системы
на независимые функциональные компоненты. Такой подход обеспечивает
упрощение сопровождения программного обеспечения, возможность расширения
функциональности системы и повышение устойчивости её работы.

В состав архитектуры приложения входят следующие основные компоненты:


В рамках рассматриваемого вопроса выделяются следующие положения: во-первых, клиентская часть приложения; во-вторых, серверная часть обработки запросов; в-третьих, подсистема интеллектуального анализа обращений; в-четвертых, подсистема хранения данных; в-пятых, подсистема передачи сформированных заявок менеджеру автомобильного сервиса.


Общая архитектура разработанной системы представлена на рисунке 2.1.

\begin{figure}[H]
\centering
\includegraphics[width=0.95\textwidth]{thesis/images/media/image10.png}
\caption{Архитектура автоматизированной системы первичной консультации клиентов автомобильного сервиса}
\label{fig:architecture}
\end{figure}

Клиентская часть приложения реализована в виде адаптивного
веб-интерфейса формата Progressive Web Application, обеспечивающего
возможность работы системы как на персональных компьютерах, так и на
мобильных устройствах без разработки отдельного нативного мобильного
приложения. Основной задачей клиентской части является организация
диалогового взаимодействия пользователя с интеллектуальным ассистентом,
ввод текстовых сообщений, отображение уточняющих вопросов и визуализация
результатов первичной консультации.

Передача пользовательских сообщений осуществляется через REST API
серверной части приложения. Клиентская часть не выполняет
интеллектуальной обработки данных и используется исключительно как
средство взаимодействия пользователя с системой.

Серверная часть приложения реализована на платформе Node.js и
обеспечивает обработку пользовательских запросов, управление
консультационными сессиями и взаимодействие с подсистемой искусственного
интеллекта. Сервер принимает текстовые сообщения пользователя, выполняет
их предварительную обработку, сохраняет историю диалога и передаёт
данные в интеллектуальный модуль анализа обращений.

Подсистема интеллектуальной обработки обращений реализована на основе
локально развёрнутой языковой модели Qwen 2.5, функционирующей через
среду выполнения Ollama. Использование локальной языковой модели
позволяет обеспечить автономность работы системы и исключить зависимость
от внешних облачных сервисов обработки данных.

Интеллектуальный модуль выполняет извлечение ключевых параметров
обращения пользователя, включая сведения об автомобиле, симптомы
неисправности и условия их проявления, формирует уточняющие вопросы при
недостаточности исходной информации, а также определяет вероятные
причины неисправности транспортного средства и формирует предварительные
диагностические рекомендации.

Результаты анализа передаются серверной части приложения, где
выполняется сопоставление полученных выводов со справочником услуг
автомобильного сервиса. На основе данного сопоставления формируется
ориентировочная стоимость обслуживания в формате «от», а также создаётся
структурированная заявка на обслуживание.

Хранение информации осуществляется с использованием системы управления
базами данных PostgreSQL. В базе данных сохраняются сведения о
пользователях системы, консультационных сессиях, сообщениях диалога,
диагностических сценариях и сформированных заявках на обслуживание.
Использование централизованного хранилища данных обеспечивает
целостность информации и возможность её последующего анализа.

Передача сформированной заявки менеджеру автомобильного сервиса
осуществляется автоматически после завершения консультационной сессии
пользователя. Это позволяет сократить время обработки обращения и
обеспечить оперативную подготовку данных для дальнейшей работы
специалистов сервиса.

Применение клиент-серверной архитектуры с интеграцией локального
интеллектуального модуля обеспечивает автономность работы системы,
масштабируемость программной структуры и возможность дальнейшего
расширения функциональности приложения за счёт добавления новых
диагностических сценариев и сервисных модулей.

\subsection{Обоснование выбора технологического
стека}\label{ux43eux431ux43eux441ux43dux43eux432ux430ux43dux438ux435-ux432ux44bux431ux43eux440ux430-ux442ux435ux445ux43dux43eux43bux43eux433ux438ux447ux435ux441ux43aux43eux433ux43e-ux441ux442ux435ux43aux430}

Разработка автоматизированной информационной системы первичной
консультации клиентов автомобильного сервиса требует использования
современного технологического стека, обеспечивающего возможность
диалогового взаимодействия с пользователем, автономной работы
интеллектуального модуля, хранения структурированных данных и
масштабируемости архитектуры приложения.

Выбор используемых технологий осуществлялся с учётом требований к
производительности системы, независимости от внешних сервисов обработки
данных, удобства сопровождения программного обеспечения и возможности
дальнейшего расширения функциональности приложения.

\subsubsection{Обоснование выбора архитектурного подхода}

Разрабатываемая система реализована на основе клиент-серверной
архитектуры, обеспечивающей разделение пользовательского интерфейса и
логики обработки данных.

Применение клиент-серверного подхода позволяет:


В рамках рассматриваемого вопроса выделяются следующие положения: во-первых, обеспечить централизованную обработку пользовательских обращений; во-вторых, организовать взаимодействие с подсистемой искусственного интеллекта через серверную часть приложения; в-третьих, реализовать хранение консультационных сессий в базе данных; в-четвертых, обеспечить возможность масштабирования системы при увеличении количества пользователей; в-пятых, упростить сопровождение и модернизацию программного обеспечения.


Использование данного архитектурного подхода является стандартным
решением при разработке интеллектуальных веб-приложений с диалоговым
интерфейсом.

\subsubsection{Обоснование выбора технологии клиентской части приложения}

Клиентская часть системы реализована в виде веб-приложения формата
Progressive Web Application (PWA), функционирующего в браузере
пользователя без необходимости установки дополнительного программного
обеспечения.

Выбор данного подхода обусловлен следующими преимуществами:


В рамках рассматриваемого вопроса выделяются следующие положения: во-первых, поддержка работы приложения на персональных компьютерах и мобильных устройствах; во-вторых, отсутствие необходимости разработки отдельных нативных мобильных приложений; в-третьих, снижение сложности сопровождения программного обеспечения; в-четвертых, обеспечение быстрого доступа пользователей к системе через стандартный веб-браузер; в-пятых, возможность последующего расширения интерфейса без изменения архитектуры системы.


Использование формата Progressive Web Application позволяет объединить
преимущества веб-приложений и мобильных решений в рамках единой
программной платформы.

\subsubsection{Обоснование выбора технологии серверной части приложения}

Серверная часть системы реализована на платформе Node.js, обеспечивающей
эффективную обработку большого количества пользовательских запросов в
асинхронном режиме.

Выбор Node.js обусловлен следующими преимуществами:


В рамках рассматриваемого вопроса выделяются следующие положения: во-первых, высокая производительность при обработке параллельных запросов; во-вторых, удобство интеграции с веб-интерфейсами; в-третьих, возможность реализации REST API для взаимодействия компонентов системы; в-четвертых, поддержка взаимодействия с библиотеками обработки естественного языка; в-пятых, широкая распространённость технологии и наличие развитой экосистемы программных модулей.


Использование платформы Node.js позволяет обеспечить стабильную работу
системы первичной консультации при увеличении количества пользователей.

\subsubsection{Обоснование выбора системы управления базами данных}

В качестве системы хранения данных используется реляционная система
управления базами данных PostgreSQL.

Выбор PostgreSQL обусловлен следующими преимуществами:


В рамках рассматриваемого вопроса выделяются следующие положения: во-первых, поддержка сложных структурированных данных; во-вторых, высокая надёжность хранения информации; в-третьих, обеспечение целостности данных; в-четвертых, возможность масштабирования базы данных; в-пятых, поддержка расширений для работы с современными типами данных; в-шестых, совместимость с серверной частью приложения.


Использование PostgreSQL обеспечивает эффективное хранение сведений о
пользователях, консультационных сессиях, диагностических сценариях и
сформированных заявках на обслуживание.

\subsubsection{Обоснование выбора технологии интеллектуальной обработки
обращений}

Ключевым компонентом разрабатываемой системы является подсистема
интеллектуального анализа пользовательских обращений, реализованная на
основе локально развёрнутой языковой модели Qwen 2.5, функционирующей
через среду выполнения Ollama.

Выбор локальной языковой модели обусловлен следующими факторами:

обеспечение автономности функционирования системы;

отсутствие зависимости от внешних облачных сервисов;


В рамках рассматриваемого вопроса выделяются следующие положения: во-первых, повышение уровня конфиденциальности пользовательских данных; во-вторых, возможность адаптации модели под предметную область автомобильного сервиса; в-третьих, снижение эксплуатационных затрат на использование внешних API.


Использование среды Ollama позволяет организовать эффективное
взаимодействие серверной части приложения с языковой моделью и
обеспечить обработку текстовых обращений пользователей в режиме
реального времени.

Применение языковой модели Qwen 2.5 обеспечивает извлечение параметров
обращения пользователя, формирование уточняющих вопросов и подготовку
предварительных диагностических рекомендаций на основе анализа
текстового описания проблемы.

\subsubsection{Обоснование выбора протокола взаимодействия компонентов
системы}

Взаимодействие между клиентской и серверной частями приложения
реализовано с использованием архитектурного стиля REST.

Использование REST API позволяет:


В рамках рассматриваемого вопроса выделяются следующие положения: во-первых, обеспечить стандартизированное взаимодействие компонентов системы; во-вторых, реализовать обмен данными в формате JSON; в-третьих, упростить интеграцию интеллектуального модуля с серверной частью приложения; в-четвертых, обеспечить масштабируемость архитектуры системы; в-пятых, повысить удобство сопровождения программного обеспечения.


Применение REST API является распространённым решением при разработке
современных веб-приложений с распределённой архитектурой.

\subsection{Разработка модели данных
приложения}\label{ux440ux430ux437ux440ux430ux431ux43eux442ux43aux430-ux43cux43eux434ux435ux43bux438-ux434ux430ux43dux43dux44bux445-ux43fux440ux438ux43bux43eux436ux435ux43dux438ux44f}

Разработка автоматизированной информационной системы первичной
консультации клиентов автомобильного сервиса требует организации
структурированного хранения данных, обеспечивающего поддержку
консультационных сессий пользователей, обработку диагностических
сценариев и формирование заявок на обслуживание.

Для хранения данных системы используется реляционная модель базы данных,
реализованная на основе системы управления базами данных PostgreSQL.
Выбор реляционной модели обусловлен необходимостью обеспечения
целостности данных, поддержки связей между сущностями системы и
возможности масштабирования структуры базы данных при расширении
функциональности приложения.

Модель данных приложения включает совокупность взаимосвязанных таблиц,
обеспечивающих хранение информации о пользователях системы,
консультационных сессиях, сообщениях диалога, диагностических сценариях
и справочных материалах автомобильного сервиса.

Основная структура базы данных автоматизированной системы представлена
на рисунке 2.3.

\begin{figure}[H]
\centering
\includegraphics[width=0.65\textwidth]{thesis/images/media/image11.png}
\caption{Модель данных автоматизированной системы первичной консультации клиентов автомобильного сервиса}
\label{fig:data-model}
\end{figure}

Описание структуры базы данных

Разработанная база данных автоматизированной системы первичной
консультации клиентов автомобильного сервиса включает совокупность
взаимосвязанных таблиц, обеспечивающих хранение информации о
пользователях системы, консультационных сессиях, диагностических
сценариях и результатах интеллектуального анализа обращений.

В состав базы данных входят следующие основные таблицы.

Таблица USERS предназначена для хранения регистрационных данных
пользователей системы. В таблице сохраняются адрес электронной почты
пользователя, хеш пароля, имя, фамилия, роль пользователя в системе и
дата регистрации.

Использование данной таблицы обеспечивает идентификацию пользователей
системы и разграничение прав доступа к функциональным возможностям
приложения.

Таблица CONSULTATION\_SESSIONS предназначена для хранения информации о
консультационных сессиях пользователей. В таблице сохраняются сведения о
пользователе, выбранном диагностическом сценарии, истории сообщений
диалога, извлечённых параметрах обращения и показателе уверенности
корректности сформированного результата.

Данные таблицы используются для формирования предварительных
диагностических выводов и последующего анализа обращений пользователей.

Таблица SERVICES содержит сведения о видах услуг автомобильного сервиса.
В таблице сохраняются наименование услуги, описание, порядок отображения
услуги в пользовательском интерфейсе и признак активности записи.

Данная таблица используется при классификации пользовательских обращений
и формировании предварительной стоимости обслуживания.

Таблица DIAGNOSTIC\_SCENARIOS предназначена для хранения диагностических
сценариев, используемых системой при анализе пользовательских обращений.
В таблице сохраняются описание возможных неисправностей, рекомендации по
их устранению, уровень критичности неисправности и условия применения
сценария.

Использование данной таблицы позволяет формировать предварительные
диагностические заключения на основе введённых пользователем симптомов.

Таблица MATERIALS содержит справочные материалы, используемые
интеллектуальной системой при обработке обращений пользователей. В
таблице сохраняются текстовые описания диагностических признаков
неисправностей и их векторные представления, применяемые для поиска
релевантной информации.

Применение данной таблицы позволяет расширить базу знаний системы и
повысить точность рекомендаций.

Таблица HINTS предназначена для хранения уточняющих вопросов,
формируемых интеллектуальным ассистентом в процессе консультации
пользователя. Использование подсказок позволяет повысить полноту
собираемой информации и улучшить качество предварительной диагностики.

\subsection{Проектирование пользовательского интерфейса
приложения}\label{ux43fux440ux43eux435ux43aux442ux438ux440ux43eux432ux430ux43dux438ux435-ux43fux43eux43bux44cux437ux43eux432ux430ux442ux435ux43bux44cux441ux43aux43eux433ux43e-ux438ux43dux442ux435ux440ux444ux435ux439ux441ux430-ux43fux440ux438ux43bux43eux436ux435ux43dux438ux44f}

Проектирование пользовательского интерфейса автоматизированной системы
первичной консультации клиентов автомобильного сервиса выполнялось с
учётом особенностей взаимодействия пользователя с интеллектуальным
ассистентом, требований к последовательности прохождения
консультационного сценария и необходимости адаптивного отображения
интерфейса на различных типах устройств.

Основной задачей разработки пользовательского интерфейса является
обеспечение удобного и интуитивно понятного взаимодействия пользователя
с системой без необходимости предварительного обучения работе с
программным обеспечением.

Интерфейс системы реализован в формате веб-приложения с поддержкой
технологии Progressive Web Application, что обеспечивает возможность
использования системы как на персональных компьютерах, так и на
мобильных устройствах.

\subsubsection{Требования к пользовательскому интерфейсу системы}

При проектировании пользовательского интерфейса учитывались следующие
требования:


В рамках рассматриваемого вопроса выделяются следующие положения: во-первых, минимизация количества действий пользователя при прохождении консультации; во-вторых, последовательное отображение этапов взаимодействия с системой; в-третьих, возможность ввода информации в свободной текстовой форме; в-четвертых, адаптивность интерфейса для различных разрешений экранов; в-пятых, обеспечение читаемости текстовой информации; в-шестых, визуальное выделение ключевых элементов управления; в-седьмых, логическая завершённость каждого этапа консультационного процесса.


Реализация указанных требований обеспечивает повышение удобства
взаимодействия пользователя с системой и снижение вероятности ошибок при
вводе данных.

\subsubsection{Структура пользовательского интерфейса системы}

Пользовательский интерфейс системы построен по сценарию
последовательного прохождения этапов консультации и включает следующие
основные экраны:

главный экран приложения;


В рамках рассматриваемого вопроса выделяются следующие положения: во-первых, экран диалогового взаимодействия с интеллектуальным ассистентом; во-вторых, экран отображения результатов предварительной диагностики; в-третьих, экран подтверждения формирования заявки на обслуживание.


Структура пользовательского интерфейса представлена на рисунке 2.4.

\begin{figure}[H]
\centering

\includegraphics[width=0.95\textwidth]{thesis/images/media/image12.png}
\caption{Структура пользовательского интерфейса системы}
\label{fig:ui-structure}
\end{figure}

\subsubsection{Проектирование главного экрана приложения}

Главный экран предназначен для первичного взаимодействия пользователя с
системой и обеспечивает быстрый доступ к запуску консультационной
сессии.

На главном экране размещаются:


В рамках рассматриваемого вопроса выделяются следующие положения: во-первых, краткое описание возможностей интеллектуального ассистента; во-вторых, элементы навигации по разделам приложения; в-третьих, кнопка запуска консультации; в-четвертых, информационные блоки с преимуществами использования системы.


Структура главного экрана представлена на рисунке 2.5.

\begin{figure}[H]
\centering
\includegraphics[width=0.95\textwidth]{thesis/images/media/image13.png}
\caption{Главный экран приложения}
\label{fig:main-screen}
\end{figure}

\subsubsection{Проектирование интерфейса диалоговой консультации}

Диалоговый интерфейс является основным элементом взаимодействия
пользователя с системой и обеспечивает ввод информации о транспортном
средстве и признаках неисправности.

В процессе консультации система последовательно запрашивает:


В рамках рассматриваемого вопроса выделяются следующие положения: во-первых, марку автомобиля; во-вторых, модель транспортного средства; в-третьих, пробег автомобиля; в-четвертых, описание симптомов неисправности.


При недостаточности исходных данных интеллектуальный ассистент формирует
уточняющие вопросы, направленные на повышение точности предварительного
диагностического заключения.

Диалог реализован в формате текстового чата, обеспечивающего
последовательное отображение сообщений пользователя и ответов системы.

Интерфейс диалоговой консультации представлен на рисунке 2.6.

\begin{figure}[H]
\centering
\includegraphics[width=0.95\textwidth]{thesis/images/media/image14.png}
\caption{Интерфейс диалоговой консультации}
\label{fig:dialog-interface}
\end{figure}

\subsubsection{Проектирование экрана отображения результатов
консультации}

После завершения анализа пользовательского обращения система формирует
предварительное диагностическое заключение и отображает результаты
обработки запроса.

Экран результатов содержит:


В рамках рассматриваемого вопроса выделяются следующие положения: во-первых, предполагаемую категорию неисправности; во-вторых, описание возможных причин возникновения проблемы; в-третьих, рекомендации по дальнейшим действиям пользователя; в-четвертых, ориентировочную стоимость обслуживания; в-пятых, показатель уверенности корректности результата.


Экран результатов консультации представлен на рисунке 2.7.

\begin{figure}[H]
\centering
\includegraphics[width=0.95\textwidth]{thesis/images/media/image15.png}
\caption{Экран отображения результатов консультации}
\label{fig:result-screen}
\end{figure}

\subsubsection{Проектирование экрана подтверждения формирования заявки}

После завершения консультационной сессии система формирует
структурированную заявку на обслуживание автомобиля и передаёт её
менеджеру автосервиса.

Пользователь получает уведомление о завершении консультации и успешной
передаче заявки.

Экран подтверждения содержит:


В рамках рассматриваемого вопроса выделяются следующие положения: во-первых, сообщение о завершении консультации; во-вторых, краткую информацию о сформированной заявке; в-третьих, уведомление о дальнейшей обработке обращения специалистом автосервиса.


Интерфейс подтверждения формирования заявки представлен на рисунке 2.8.

\begin{figure}[H]
\centering
\includegraphics[width=0.95\textwidth]{thesis/images/media/image16.png}
\caption{Экран подтверждения формирования заявки}
\label{fig:confirm-screen}
\end{figure}

\subsection{Реализация серверной части
системы}\label{ux440ux435ux430ux43bux438ux437ux430ux446ux438ux44f-ux441ux435ux440ux432ux435ux440ux43dux43eux439-ux447ux430ux441ux442ux438-ux441ux438ux441ux442ux435ux43cux44b}

Серверная часть автоматизированной системы первичной консультации
клиентов автомобильного сервиса обеспечивает обработку пользовательских
обращений, управление консультационными сессиями, взаимодействие с
подсистемой искусственного интеллекта и формирование структурированных
заявок на обслуживание.

Реализация серверной части выполнена на платформе Node.js с
использованием фреймворка Express, обеспечивающего обработку
HTTP-запросов и организацию REST API приложения. Использование
асинхронной модели обработки запросов позволяет обеспечить стабильную
работу системы при одновременном взаимодействии нескольких
пользователей.

Архитектура серверного приложения построена по модульному принципу, что
обеспечивает масштабируемость системы и упрощает сопровождение
программного обеспечения.

\subsubsection{Общая структура серверной части приложения}

Серверная часть системы включает следующие функциональные компоненты:


В рамках рассматриваемого вопроса выделяются следующие положения: во-первых, модуль обработки HTTP-запросов пользователя; во-вторых, модуль маршрутизации REST API; в-третьих, модуль управления консультационными сессиями; в-четвертых, модуль взаимодействия с базой данных; в-пятых, модуль интеграции с языковой моделью; в-шестых, модуль формирования заявок на обслуживание.


Инициализация серверного приложения включает настройку политики
безопасности, конфигурацию CORS, подключение middleware-обработчиков и
регистрацию маршрутов REST API системы.

Пример реализации инициализации серверного приложения представлен в
\hyperref[ux43bux438ux441211]{листинге 2.1}.

{\def\LTcaptype{none} % do not increment counter
\begin{longtable}[]{@{}
  >{\raggedright\arraybackslash}p{(\linewidth - 2\tabcolsep) * \real{0.0938}}
  >{\raggedright\arraybackslash}p{(\linewidth - 2\tabcolsep) * \real{0.9062}}@{}}
\toprule\noalign{}
\begin{minipage}[b]{\linewidth}\raggedright
1.

2.

3.

4.

5.

6.

7.

8.

9.

10.

11.

12.

13.

14.

15.

16.

17.

18.

19.

20.

21.

22.

23.

24.

25.

26.

27.
\end{minipage} & \begin{minipage}[b]{\linewidth}\raggedright
export function createApp() \{

const app = express();

const env = getEnv();

if (env.NODE\_ENV === \textquotesingle production\textquotesingle) \{

app.set(\textquotesingle trust proxy\textquotesingle, 1);

\}

app.use(

helmet(\{

contentSecurityPolicy: false,

\}),

);

app.use(

cors(\{

origin: true,

credentials: true,
\end{minipage} \\
\midrule\noalign{}
\endhead
\bottomrule\noalign{}
\endlastfoot
\end{longtable}
}

Продолжение листинга 2.1

{\def\LTcaptype{none} % do not increment counter
\begin{longtable}[]{@{}
  >{\raggedright\arraybackslash}p{(\linewidth - 2\tabcolsep) * \real{0.0938}}
  >{\raggedright\arraybackslash}p{(\linewidth - 2\tabcolsep) * \real{0.9062}}@{}}
\toprule\noalign{}
\begin{minipage}[b]{\linewidth}\raggedright
24.

25.

26.

27.

28.

29.

30.

31.

32.

33.

34.

35.

36.

37.

38.

39.

40.

41.

42.

43.

44.

45.

46.

47.

48.
\end{minipage} & \begin{minipage}[b]{\linewidth}\raggedright
\}),

);

app.use(cookieParser());

app.use(express.json(\{ limit: \textquotesingle1mb\textquotesingle{}
\}));

const limiter = rateLimit(\{

windowMs: 15 * 60 * 1000,

max: 300,

\});

app.use(\textquotesingle/api/\textquotesingle, limiter);

app.use(\textquotesingle/api\textquotesingle, api);

return app;

\}
\end{minipage} \\
\midrule\noalign{}
\endhead
\bottomrule\noalign{}
\endlastfoot
\end{longtable}
}

\hyperref[ux43bux438ux441211]{Листинг 2.1} -- Инициализация серверного
приложения Express

Представленный фрагмент демонстрирует подключение middleware обеспечения
безопасности, ограничение частоты запросов пользователей и регистрацию
маршрутов REST API системы.

\subsubsection{Реализация REST API серверной части}

Взаимодействие клиентской и серверной частей приложения реализовано с
использованием архитектурного стиля REST. REST API обеспечивает передачу
пользовательских сообщений, получение результатов интеллектуального
анализа и управление консультационными сессиями.

Серверная часть поддерживает следующие группы маршрутов:


В рамках рассматриваемого вопроса выделяются следующие положения: во-первых, маршруты аутентификации пользователей; во-вторых, маршруты обработки консультационных сессий; в-третьих, маршруты формирования заявок на обслуживание; в-четвертых, маршруты административного управления системой; в-пятых, маршруты аналитики обращений пользователей.


Для контроля работоспособности серверной части реализован специальный
маршрут проверки состояния базы данных. Пример реализации данного
маршрута представлен в
\protect\phantomsection\label{ux43bux438ux44122}{}\hyperref[ux43bux438ux441221]{листинге
2.2.}

{\def\LTcaptype{none} % do not increment counter
\begin{longtable}[]{@{}
  >{\raggedright\arraybackslash}p{(\linewidth - 2\tabcolsep) * \real{0.0938}}
  >{\raggedright\arraybackslash}p{(\linewidth - 2\tabcolsep) * \real{0.9062}}@{}}
\toprule\noalign{}
\begin{minipage}[b]{\linewidth}\raggedright
1.

2.

3.

4.

5.

6.

7.

8.

9.

10.

11.

12.

13.
\end{minipage} & \begin{minipage}[b]{\linewidth}\raggedright
api.get(\textquotesingle/health\textquotesingle, async (\_req, res)
=\textgreater{} \{

try \{

await prisma.\$queryRaw`SELECT 1`;

res.json(\{ status: \textquotesingle ok\textquotesingle, db:
\textquotesingle connected\textquotesingle{} \});

\} catch (err) \{

res.status(503).json(\{ status: \textquotesingle error\textquotesingle,
db: \textquotesingle disconnected\textquotesingle{} \});

\}

\});
\end{minipage} \\
\midrule\noalign{}
\endhead
\bottomrule\noalign{}
\endlastfoot
\end{longtable}
}

\hyperref[ux43bux438ux44122]{Листинг 2.2} -- Проверка доступности базы
данных через REST API

Использование данного маршрута позволяет контролировать доступность базы
данных и своевременно выявлять возможные ошибки функционирования
серверной части системы.

\subsubsection{Управление консультационными сессиями пользователей}

Модуль управления консультационными сессиями обеспечивает сохранение
истории взаимодействия пользователя с интеллектуальным ассистентом и
последовательную обработку пользовательских сообщений.

В рамках консультационной сессии выполняются следующие операции:


В рамках рассматриваемого вопроса выделяются следующие положения: во-первых, создание новой консультационной сессии; во-вторых, сохранение сообщений пользователя; в-третьих, сохранение ответов интеллектуального ассистента; в-четвертых, накопление извлечённых параметров обращения; в-пятых, расчёт показателя уверенности корректности результата; в-шестых, завершение консультационной сессии после формирования заявки.


Использование механизма консультационных сессий обеспечивает сохранение
контекста диалога и повышение точности предварительной диагностики.

\subsubsection{Интеграция серверной части с подсистемой искусственного
интеллекта}

Серверная часть системы обеспечивает взаимодействие с локально
развёрнутой языковой моделью через среду выполнения Ollama.

После получения текстового сообщения пользователя выполняется передача
данных в интеллектуальный модуль, где осуществляется:


В рамках рассматриваемого вопроса выделяются следующие положения: во-первых, анализ текстового описания проблемы; во-вторых, извлечение параметров обращения пользователя; в-третьих, формирование уточняющих вопросов; в-четвертых, определение категории неисправности; в-пятых, формирование диагностических рекомендаций; в-шестых, расчёт показателя уверенности корректности результата.


Полученные результаты передаются серверной части приложения и
используются при формировании структурированной заявки на обслуживание.

Использование локальной языковой модели обеспечивает автономность
функционирования системы и повышает уровень конфиденциальности
пользовательских данных.

\subsubsection{Взаимодействие серверной части с базой данных}

Серверная часть приложения обеспечивает сохранение и обработку
информации пользователей, консультационных сессий и диагностических
сценариев в базе данных PostgreSQL.

В рамках взаимодействия с базой данных выполняются следующие операции:


В рамках рассматриваемого вопроса выделяются следующие положения: во-первых, сохранение регистрационных данных пользователей; во-вторых, сохранение сообщений консультационных сессий; в-третьих, сохранение извлечённых параметров обращения; в-четвертых, получение диагностических сценариев; в-пятых, формирование и сохранение заявок на обслуживание; в-шестых, сохранение уточняющих подсказок интеллектуального ассистента.


Использование централизованного хранилища данных обеспечивает
целостность информации и возможность её последующего анализа.

\subsubsection{Формирование структурированной заявки на обслуживание}

После завершения обработки пользовательского обращения серверная часть
системы формирует структурированную заявку на обслуживание автомобиля.

В состав заявки включаются:


В рамках рассматриваемого вопроса выделяются следующие положения: во-первых, сведения о пользователе; во-вторых, параметры транспортного средства; в-третьих, описание симптомов неисправности; в-четвертых, результаты предварительной диагностики; в-пятых, показатель уверенности корректности результата; в-шестых, ориентировочная стоимость обслуживания.


Сформированная заявка сохраняется в базе данных и передаётся менеджеру
автомобильного сервиса для дальнейшей обработки.

Автоматизация формирования заявок позволяет сократить время обработки
обращений клиентов и повысить эффективность работы специалистов сервиса.

\subsection{Реализация клиентской части
системы}\label{ux440ux435ux430ux43bux438ux437ux430ux446ux438ux44f-ux43aux43bux438ux435ux43dux442ux441ux43aux43eux439-ux447ux430ux441ux442ux438-ux441ux438ux441ux442ux435ux43cux44b}

Клиентская часть автоматизированной системы первичной консультации
клиентов автомобильного сервиса обеспечивает взаимодействие пользователя
с интеллектуальным ассистентом и реализует интерфейс прохождения
консультационной сессии.

Клиентская часть разработана в формате веб-приложения с поддержкой
технологии Progressive Web Application (PWA), что обеспечивает
возможность использования системы как на персональных компьютерах, так и
на мобильных устройствах без установки дополнительного программного
обеспечения.

Основными задачами клиентской части являются:


В рамках рассматриваемого вопроса выделяются следующие положения: во-первых, отображение пользовательского интерфейса системы; во-вторых, организация диалогового взаимодействия пользователя с интеллектуальным ассистентом; в-третьих, передача пользовательских сообщений серверной части приложения; в-четвертых, получение результатов интеллектуального анализа обращения; в-пятых, отображение результатов предварительной диагностики; в-шестых, информирование пользователя о формировании заявки на обслуживание.


Архитектура клиентской части построена по модульному принципу, что
обеспечивает удобство сопровождения программного обеспечения и
возможность расширения функциональности интерфейса.

\subsubsection{Архитектура клиентского приложения}

Клиентское приложение реализовано в виде набора взаимосвязанных
программных модулей, обеспечивающих взаимодействие интерфейса
пользователя с серверной частью системы.

В состав клиентской части входят следующие основные модули:


В рамках рассматриваемого вопроса выделяются следующие положения: во-первых, модуль отображения пользовательского интерфейса; во-вторых, модуль взаимодействия с REST API; в-третьих, модуль управления консультационной сессией; в-четвертых, модуль отображения сообщений диалогового интерфейса; в-пятых, модуль обработки ошибок пользовательских запросов.


Использование модульной архитектуры позволяет обеспечить независимость
компонентов пользовательского интерфейса и упростить их сопровождение.

\subsubsection{Взаимодействие клиентской части с сервером}

Взаимодействие клиентской части приложения с серверной частью
реализовано через REST API с использованием универсального программного
интерфейса обмена данными.

Разработанный API-клиент обеспечивает:


В рамках рассматриваемого вопроса выделяются следующие положения: во-первых, автоматическое добавление заголовков запросов; во-вторых, поддержку CSRF-защиты; в-третьих, обработку ошибок сервера; в-четвертых, повторную отправку запросов при обновлении пользовательской сессии; в-пятых, преобразование полученных данных в структурированный формат.


Пример реализации универсального API-клиента представлен в
\protect\phantomsection\label{ux43bux438ux44123}{}\hyperref[ux43bux438ux441231]{листинге
2.3.}

{\def\LTcaptype{none} % do not increment counter
\begin{longtable}[]{@{}
  >{\raggedright\arraybackslash}p{(\linewidth - 2\tabcolsep) * \real{0.0938}}
  >{\raggedright\arraybackslash}p{(\linewidth - 2\tabcolsep) * \real{0.9062}}@{}}
\toprule\noalign{}
\begin{minipage}[b]{\linewidth}\raggedright
1.

2.

3.

4.

5.

6.

7.

8.

9.

10.
\end{minipage} & \begin{minipage}[b]{\linewidth}\raggedright
async function rawFetch(path, opts = \{\}) \{

const headers = \{ ..(opts.headers \textbar\textbar{} \{\}) \};

if (opts.body != null \&\&

!(opts.body instanceof FormData) \&\&

!headers{[}\textquotesingle Content-Type\textquotesingle{]}) \{

headers{[}\textquotesingle Content-Type\textquotesingle{]} =
\textquotesingle application/json\textquotesingle;

\}

if (mutatingMethod(opts.method) \&\& !opts.skipCsrf) \{
\end{minipage} \\
\midrule\noalign{}
\endhead
\bottomrule\noalign{}
\endlastfoot
\end{longtable}
}

Продолжение листинга 2.3

{\def\LTcaptype{none} % do not increment counter
\begin{longtable}[]{@{}
  >{\raggedright\arraybackslash}p{(\linewidth - 2\tabcolsep) * \real{0.0938}}
  >{\raggedright\arraybackslash}p{(\linewidth - 2\tabcolsep) * \real{0.9062}}@{}}
\toprule\noalign{}
\begin{minipage}[b]{\linewidth}\raggedright
4.

5.

6.

7.

8.

9.

10.

11.

12.

13.

14.

15.

16.

17.

18.

19.

20.

21.

22.

23.

24.

25.

26.

27.

28.

29.

30.

31.

32.

33.

34.

35.

36.

37.
\end{minipage} & \begin{minipage}[b]{\linewidth}\raggedright
if (opts.body != null \&\&

!(opts.body instanceof FormData) \&\&

!headers{[}\textquotesingle Content-Type\textquotesingle{]}) \{

headers{[}\textquotesingle Content-Type\textquotesingle{]} =
\textquotesingle application/json\textquotesingle;

\}

if (mutatingMethod(opts.method) \&\& !opts.skipCsrf) \{

const csrf = readCookie(\textquotesingle fm\_csrf\textquotesingle);

if (csrf) headers{[}\textquotesingle X-CSRF-Token\textquotesingle{]} =
csrf;

\}

const res = await fetch(`\$\{API\_BASE\}\$\{path\}`, \{

..opts,

credentials: \textquotesingle include\textquotesingle,

headers,

body:

opts.body != null \&\&

typeof opts.body === \textquotesingle object\textquotesingle{} \&\&

!(opts.body instanceof FormData)

? JSON.stringify(opts.body)

: opts.body,

\});

return res;

\} if (opts.body != null \&\&

!(opts.body instanceof FormData) \&\&

!headers{[}\textquotesingle Content-Type\textquotesingle{]}) \{

headers{[}\textquotesingle Content-Type\textquotesingle{]} =
\textquotesingle application/json\textquotesingle;

\}

if (mutatingMethod(opts.method) \&\& !opts.skipCsrf) \{

const csrf = readCookie(\textquotesingle fm\_csrf\textquotesingle);

if (csrf) headers{[}\textquotesingle X-CSRF-Token\textquotesingle{]} =
csrf;

\}
\end{minipage} \\
\midrule\noalign{}
\endhead
\bottomrule\noalign{}
\endlastfoot
\end{longtable}
}

\hyperref[ux43bux438ux44123]{Листинг 2.3} -- Универсальный API-клиент
клиентской части приложения

Использование универсального API-клиента позволяет централизованно
управлять взаимодействием пользовательского интерфейса с серверной
частью системы и обеспечивает повышение устойчивости работы приложения.

\subsubsection{Реализация диалогового интерфейса консультации}

Основным элементом взаимодействия пользователя с системой является
диалоговый интерфейс консультации, реализованный в виде текстового чата.

Передача сообщений пользователя интеллектуальному ассистенту
осуществляется с использованием потоковой передачи данных по протоколу
Server-Sent Events (SSE). Использование данного подхода позволяет
отображать ответы системы по мере их формирования, что повышает удобство
взаимодействия пользователя с приложением.

Использование потоковой передачи сообщений позволяет обеспечить
отображение ответов интеллектуального ассистента в режиме реального
времени и повысить удобство взаимодействия пользователя с системой.

\subsubsection{Отображение результатов консультации пользователю}

После завершения обработки пользовательского обращения клиентская часть
приложения получает результаты интеллектуального анализа от серверной
части системы и отображает их пользователю в структурированном виде.

На экране результатов консультации отображаются:


В рамках рассматриваемого вопроса выделяются следующие положения: во-первых, предполагаемая категория неисправности; во-вторых, возможные причины возникновения проблемы; в-третьих, рекомендации по дальнейшим действиям; в-четвертых, ориентировочная стоимость обслуживания; в-пятых, показатель уверенности корректности результата.


Отображение результатов предварительной диагностики позволяет
пользователю получить первичную оценку состояния транспортного средства
до обращения в сервисный центр.

\subsubsection{Обработка ошибок пользовательских запросов}

Для обеспечения устойчивости работы пользовательского интерфейса
реализован механизм обработки ошибок взаимодействия с серверной частью
системы.

В случае возникновения ошибок выполняются следующие действия:

отображение уведомления пользователю;


В рамках рассматриваемого вопроса выделяются следующие положения: во-первых, повторная отправка запроса при восстановлении соединения; во-вторых, автоматическое обновление пользовательской сессии; в-третьих, восстановление состояния консультационной сессии.


Использование механизма обработки ошибок позволяет повысить надёжность
работы клиентской части приложения и обеспечить корректность
взаимодействия пользователя с системой.

\subsection{Интеграция подсистемы искусственного
интеллекта}\label{ux438ux43dux442ux435ux433ux440ux430ux446ux438ux44f-ux43fux43eux434ux441ux438ux441ux442ux435ux43cux44b-ux438ux441ux43aux443ux441ux441ux442ux432ux435ux43dux43dux43eux433ux43e-ux438ux43dux442ux435ux43bux43bux435ux43aux442ux430}

Подсистема искусственного интеллекта является ключевым компонентом
автоматизированной системы первичной консультации клиентов
автомобильного сервиса и обеспечивает анализ текстовых обращений
пользователей, извлечение диагностически значимых параметров и
формирование предварительных рекомендаций по обслуживанию транспортного
средства.

Интеллектуальная обработка обращений реализована на основе локально
развёрнутой языковой модели, взаимодействующей с серверной частью
приложения через специализированный программный интерфейс среды
выполнения Ollama. Использование локальной модели позволяет обеспечить
автономность работы системы и повысить уровень конфиденциальности
пользовательских данных.

Архитектура подсистемы искусственного интеллекта включает следующие
основные компоненты:


В рамках рассматриваемого вопроса выделяются следующие положения: во-первых, модуль формирования запросов к языковой модели; во-вторых, модуль предварительного анализа симптомов; в-третьих, модуль объединения результатов rule-based анализа и вывода модели; в-четвертых, модуль структурирования диагностического результата; в-пятых, модуль расчёта показателя уверенности корректности результата.


Совместная работа указанных компонентов обеспечивает формирование
диагностического заключения на основе текстового описания проблемы
пользователем.

\subsubsection{Взаимодействие серверной части системы с языковой моделью}

Взаимодействие серверной части приложения с языковой моделью
осуществляется через HTTP API среды выполнения Ollama. Передача данных
выполняется в формате JSON, содержащем историю сообщений
консультационной сессии и параметры запроса диагностики.

Пример реализации функции взаимодействия с языковой моделью представлен
в
\protect\phantomsection\label{ux43bux438ux44124}{}\hyperref[ux43bux438ux441241]{листинге
2.4.}

{\def\LTcaptype{none} % do not increment counter
\begin{longtable}[]{@{}
  >{\raggedright\arraybackslash}p{(\linewidth - 2\tabcolsep) * \real{0.0943}}
  >{\raggedright\arraybackslash}p{(\linewidth - 2\tabcolsep) * \real{0.8748}}@{}}
\toprule\noalign{}
\begin{minipage}[b]{\linewidth}\raggedright
1.

2.

3.

4.

5.

6.

7.

8.

9.

10.

11.

12.

13.

14.

15.

16.

17.

18.

19.

20.

21.

22.

23.

24.

25.

26.

27.

28.
\end{minipage} & \begin{minipage}[b]{\linewidth}\raggedright
export async function chatCompletion(\{ messages, model, temperature =
0, format \}) \{

const env = getEnv();

if (!env.LLM\_ENABLED) \{

throw new AppError(503, \textquotesingle LLM disabled via
configuration\textquotesingle,
\textquotesingle LLM\_ERROR\textquotesingle);

\}

const base = resolveBaseUrl(env.LLM\_BASE\_URL);

const url = `\$\{base\}/api/chat`;

const body = \{

model: model \textbar\textbar{} env.LLM\_MODEL,

messages,

temperature,

stream: false,

\};

if (format) body.format = format;

const res = await fetch(url, \{

method: \textquotesingle POST\textquotesingle,
\end{minipage} \\
\midrule\noalign{}
\endhead
\bottomrule\noalign{}
\endlastfoot
\end{longtable}
}

Продолжение листинга 2.4

{\def\LTcaptype{none} % do not increment counter
\begin{longtable}[]{@{}
  >{\raggedright\arraybackslash}p{(\linewidth - 2\tabcolsep) * \real{0.0943}}
  >{\raggedright\arraybackslash}p{(\linewidth - 2\tabcolsep) * \real{0.8748}}@{}}
\toprule\noalign{}
\begin{minipage}[b]{\linewidth}\raggedright
29.

30.

31.

32.

33.

34.

35.

36.

37.
\end{minipage} & \begin{minipage}[b]{\linewidth}\raggedright
headers: \{ \textquotesingle Content-Type\textquotesingle:
\textquotesingle application/json\textquotesingle{} \},

body: JSON.stringify(body),

signal: AbortSignal.timeout(120\_000),

\});

const data = await res.json();

return String(data?.message?.content);

\}
\end{minipage} \\
\midrule\noalign{}
\endhead
\bottomrule\noalign{}
\endlastfoot
\end{longtable}
}

\hyperref[ux43bux438ux44124]{Листинг 2.4} -- Отправка запроса к языковой
модели через Ollama API

Использование локального API языковой модели позволяет исключить
зависимость системы от внешних облачных сервисов и обеспечить обработку
пользовательских обращений в защищённой среде выполнения

\subsubsection{Формирование диагностического заключения на основе
гибридного алгоритма анализа}

Формирование диагностического заключения выполняется на основе
гибридного алгоритма, объединяющего методы предварительного rule-based
анализа симптомов и интеллектуальной обработки данных языковой моделью.

На первом этапе выполняется проверка полноты исходных параметров
обращения пользователя. При недостаточности информации используется
резервный сценарий формирования диагностического заключения.

После этого формируется структурированный набор диагностических
параметров, включающий сведения о транспортном средстве, пробеге,
симптомах неисправности и условиях их проявления. Полученные данные
передаются в языковую модель для формирования предварительного
диагностического результата.

Пример реализации алгоритма формирования диагностического заключения
представлен в
\protect\phantomsection\label{ux43bux438ux44125}{}\hyperref[ux43bux438ux441251]{листинге
2.5.}

{\def\LTcaptype{none} % do not increment counter
\begin{longtable}[]{@{}
  >{\raggedright\arraybackslash}p{(\linewidth - 2\tabcolsep) * \real{0.0943}}
  >{\raggedright\arraybackslash}p{(\linewidth - 2\tabcolsep) * \real{0.8748}}@{}}
\toprule\noalign{}
\begin{minipage}[b]{\linewidth}\raggedright
1.

2.

3.

4.

5.

6.

7.

8.

9.

10.

11.

12.

13.

14.

15.

16.

17.

18.

19.

20.

21.

22.

23.

24.

25.

26.

27.

28.

29.

30.

31.

32.

33.

34.

35.

36.

37.

38.

39.

40.

41.

42.

43.

44.

45.

46.

47.
\end{minipage} & \begin{minipage}[b]{\linewidth}\raggedright
export async function generateDiagnosis(data) \{

const cond = data?.conditions ?? data?.problemConditions;

if (

!isFieldFilled(\textquotesingle mileage\textquotesingle, data.mileage)
\textbar\textbar{}

!isFieldFilled(\textquotesingle symptoms\textquotesingle, data.symptoms)
\textbar\textbar{}

!isFieldFilled(\textquotesingle conditions\textquotesingle, cond)

) \{

return safeDiagnosisFallback();

\}

const payload = \{

car\_make: data.car\_make ?? null,

car\_model: data.car\_model ?? null,

year: data.year ?? null,

mileage: data.mileage ?? null,

symptoms: data.symptoms ?? null,

conditions: cond ?? null,

urgency\_signs: data.urgency\_signs ?? null,

category: data.category ?? null,

\};

const ruleBased = preAnalyzeSymptoms(payload);

try \{

const raw = await chatCompletion(\{

temperature: 0.2,

format: DIAGNOSIS\_FORMAT\_SCHEMA,

messages: {[}

\{ role: \textquotesingle system\textquotesingle, content:
DIAGNOSIS\_SYSTEM\_PROMPT \},

\{ role: \textquotesingle user\textquotesingle, content:
diagnosisUserPrompt(payload) \},

{]},
\end{minipage} \\
\midrule\noalign{}
\endhead
\bottomrule\noalign{}
\endlastfoot
\end{longtable}
}

Продолжение листинга 2.5

{\def\LTcaptype{none} % do not increment counter
\begin{longtable}[]{@{}
  >{\raggedright\arraybackslash}p{(\linewidth - 2\tabcolsep) * \real{0.0943}}
  >{\raggedright\arraybackslash}p{(\linewidth - 2\tabcolsep) * \real{0.8748}}@{}}
\toprule\noalign{}
\begin{minipage}[b]{\linewidth}\raggedright
49.

50.

51.

52.

53.

54.

55.

56.

57.

58.

59.

60.

61.

62.

63.

64.

65.

66.

67.

68.

69.

70.

71.

72.

73.
\end{minipage} & \begin{minipage}[b]{\linewidth}\raggedright
const raw = await chatCompletion(\{

temperature: 0.2,

format: DIAGNOSIS\_FORMAT\_SCHEMA,

messages: {[}

\{ role: \textquotesingle system\textquotesingle, content:
DIAGNOSIS\_SYSTEM\_PROMPT \},

\{ role: \textquotesingle user\textquotesingle, content:
diagnosisUserPrompt(payload) \},

{]},

\});

const llmDiagnosis = normalizeDiagnosis(JSON.parse(raw));

return mergeDiagnosis(ruleBased, llmDiagnosis);

\} catch \{

return mergeDiagnosis(ruleBased, safeDiagnosisFallback());

\}

\}
\end{minipage} \\
\midrule\noalign{}
\endhead
\bottomrule\noalign{}
\endlastfoot
\end{longtable}
}

\hyperref[ux43bux438ux44125]{Листинг 2.5} -- Алгоритм формирования
диагностического заключения

Использование гибридного алгоритма позволяет повысить устойчивость
работы системы и обеспечить корректное формирование диагностического
результата даже при неполноте исходных данных пользователя.

\subsubsection{Формирование структурированного ответа языковой модели}

Для повышения надёжности обработки результатов интеллектуального анализа
используется механизм структурирования ответа языковой модели на основе
JSON-схемы.

Применение структурированного формата ответа позволяет обеспечить
корректность интерпретации результатов диагностики серверной частью
приложения и упростить дальнейшую обработку данных.

Пример схемы структурированного ответа представлен в
\protect\phantomsection\label{ux43bux438ux44126}{}\hyperref[ux43bux438ux441261]{листинге
2.6.}

{\def\LTcaptype{none} % do not increment counter
\begin{longtable}[]{@{}
  >{\raggedright\arraybackslash}p{(\linewidth - 2\tabcolsep) * \real{0.0943}}
  >{\raggedright\arraybackslash}p{(\linewidth - 2\tabcolsep) * \real{0.8748}}@{}}
\toprule\noalign{}
\begin{minipage}[b]{\linewidth}\raggedright
1.

2.

3.

4.

5.

6.

7.

8.

9.

10.

11.

12.

13.

14.

15.

16.

17.

18.

19.

20.

21.

22.

23.

24.

25.

26.

27.
\end{minipage} & \begin{minipage}[b]{\linewidth}\raggedright
export const DIAGNOSIS\_FORMAT\_SCHEMA = \{

type: \textquotesingle object\textquotesingle,

properties: \{

probable\_causes: \{ type: \textquotesingle array\textquotesingle,
items: \{ type: \textquotesingle string\textquotesingle{} \} \},

recommended\_checks: \{ type: \textquotesingle array\textquotesingle,
items: \{ type: \textquotesingle string\textquotesingle{} \} \},

urgency: \{ type: \textquotesingle string\textquotesingle, enum:
{[}\textquotesingle low\textquotesingle,
\textquotesingle medium\textquotesingle,
\textquotesingle high\textquotesingle{]} \},

confidence: \{ type: \textquotesingle number\textquotesingle{} \},

estimated\_cost\_from: \{ type:
{[}\textquotesingle integer\textquotesingle,
\textquotesingle null\textquotesingle{]} \},

summary: \{ type: \textquotesingle string\textquotesingle{} \},

\},

required: {[}

\textquotesingle probable\_causes\textquotesingle,

\textquotesingle recommended\_checks\textquotesingle,

\textquotesingle urgency\textquotesingle,

\textquotesingle confidence\textquotesingle,

\textquotesingle estimated\_cost\_from\textquotesingle,

\textquotesingle summary\textquotesingle,

{]},

\};
\end{minipage} \\
\midrule\noalign{}
\endhead
\bottomrule\noalign{}
\endlastfoot
\end{longtable}
}

\hyperref[ux43bux438ux44126]{Листинг 2.6} -- JSON-схема
структурированного ответа языковой модели

Использование формализованной схемы ответа обеспечивает стандартизацию
результатов интеллектуального анализа и повышает точность формирования
предварительных рекомендаций по обслуживанию транспортного средства.

\subsection{Тестирование системы по критическому пути
пользователя}\label{ux442ux435ux441ux442ux438ux440ux43eux432ux430ux43dux438ux435-ux441ux438ux441ux442ux435ux43cux44b-ux43fux43e-ux43aux440ux438ux442ux438ux447ux435ux441ux43aux43eux43cux443-ux43fux443ux442ux438-ux43fux43eux43bux44cux437ux43eux432ux430ux442ux435ux43bux44f}

Для проверки корректности функционирования разработанной
автоматизированной системы первичной консультации клиентов
автомобильного сервиса выполнено тестирование по критическому пути
пользователя.

Критический путь пользователя представляет собой последовательность
действий клиента автосервиса от начала взаимодействия с системой до
получения результатов предварительной диагностики неисправности
транспортного средства и формирования рекомендаций по дальнейшему
обслуживанию.

В качестве основной роли рассматривается пользователь автомобильного
сервиса, обращающийся за первичной консультацией по поводу технического
состояния автомобиля.

Цель тестирования заключается в подтверждении корректности выполнения
системой основных функциональных требований при прохождении полного
сценария консультации.

Перед началом тестирования пользователь открывает веб-приложение системы
и переходит к разделу интеллектуальной консультации.

Сценарий тестирования представлен в таблице 2.1.

{\def\LTcaptype{none} % do not increment counter
\begin{longtable}[]{@{}
  >{\centering\arraybackslash}p{(\linewidth - 6\tabcolsep) * \real{0.2500}}
  >{\centering\arraybackslash}p{(\linewidth - 6\tabcolsep) * \real{0.2500}}
  >{\centering\arraybackslash}p{(\linewidth - 6\tabcolsep) * \real{0.2500}}
  >{\centering\arraybackslash}p{(\linewidth - 6\tabcolsep) * \real{0.2500}}@{}}
\toprule\noalign{}
\begin{minipage}[b]{\linewidth}\centering
\textbf{№}
\end{minipage} & \begin{minipage}[b]{\linewidth}\centering
\textbf{Действие пользователя}
\end{minipage} & \begin{minipage}[b]{\linewidth}\centering
\textbf{Ожидаемый результат}
\end{minipage} & \begin{minipage}[b]{\linewidth}\centering
\textbf{Фактический результат}
\end{minipage} \\
\midrule\noalign{}
\endhead
\bottomrule\noalign{}
\endlastfoot
1 & Открытие страницы интеллектуальной консультации & Отображается
интерфейс диалоговой системы & Выполнено \\
2 & Ввод марки, модели и пробега автомобиля & Система принимает и
сохраняет параметры автомобиля & Выполнено \\
3 & Ввод описания симптомов неисправности & Сообщение передается
интеллектуальному модулю обработки & Выполнено \\
4 & Анализ введённых данных системой & Выполняется извлечение
диагностически значимых параметров обращения & Выполнено \\
5 & Получение уточняющих вопросов от системы & Отображаются
дополнительные вопросы для уточнения симптомов & Выполнено \\
6 & Ввод уточняющей информации пользователем & Система дополняет
диагностические параметры обращения & Выполнено \\
7 & Формирование предварительного диагностического результата &
Отображаются вероятные причины неисправности & Выполнено \\
8 & Отображение рекомендаций по проверке автомобиля & Пользователь
получает список рекомендуемых диагностических действий & Выполнено \\
9 & Отображение ориентировочной стоимости обслуживания & Отображается
предварительная стоимость ремонта & Выполнено \\
\end{longtable}
}

\subsection{Выводы по
главе}\label{ux432ux44bux432ux43eux434ux44b-ux43fux43e-ux433ux43bux430ux432ux435-1}

Во второй главе выпускной квалификационной работы рассмотрены вопросы
технической реализации автоматизированной системы первичной консультации
клиентов автомобильного сервиса на основе технологий искусственного
интеллекта.

В ходе выполнения главы разработана архитектура программного решения,
обеспечивающая взаимодействие клиентской части приложения, серверной
части системы, базы данных и подсистемы интеллектуального анализа
пользовательских обращений. Реализованная архитектура построена по
клиент-серверной модели и обеспечивает масштабируемость, модульность и
возможность дальнейшего расширения функциональности системы.

Определён технологический стек разработки, включающий платформу Node.js
для реализации серверной части приложения, систему управления базами
данных PostgreSQL для хранения структурированных данных пользователей и
консультационных сессий, а также локально развёрнутую языковую модель,
функционирующую через среду выполнения Ollama. Выбранный стек
обеспечивает автономность работы системы и исключает зависимость от
внешних облачных сервисов.

Разработана модель данных приложения, включающая сущности пользователей,
консультационных сессий, диагностических сценариев, подсказок и
справочника услуг. Построенная структура базы данных обеспечивает
хранение истории взаимодействия пользователей с системой и формирование
структурированных заявок на обслуживание транспортных средств.

Спроектирован пользовательский интерфейс системы, реализующий диалоговое
взаимодействие клиента с интеллектуальным ассистентом. Разработанные
интерфейсные решения обеспечивают последовательный сбор диагностически
значимых параметров обращения пользователя, включая сведения о
транспортном средстве, пробеге и симптомах неисправности.

Реализована серверная часть системы, обеспечивающая обработку
пользовательских запросов, управление консультационными сессиями,
взаимодействие с базой данных и формирование предварительных
диагностических заключений. Организация взаимодействия компонентов
системы выполнена с использованием архитектурного стиля REST.

Разработана клиентская часть приложения, обеспечивающая передачу
пользовательских сообщений серверной части системы и отображение
результатов предварительной диагностики в режиме реального времени.
Использование технологии потоковой передачи данных позволило повысить
удобство взаимодействия пользователя с интеллектуальным ассистентом.

Реализована интеграция подсистемы искусственного интеллекта,
обеспечивающая анализ текстовых обращений пользователей на основе
гибридного алгоритма, объединяющего rule-based методы и обработку данных
языковой моделью. Использование структурированного формата ответа модели
позволило обеспечить корректную интерпретацию результатов
интеллектуального анализа и повысить точность формирования
диагностических рекомендаций.

Таким образом, во второй главе выполнена техническая реализация
автоматизированной системы первичной консультации клиентов
автомобильного сервиса, обеспечивающая диалоговое взаимодействие
пользователя с интеллектуальным ассистентом, предварительную диагностику
неисправностей транспортного средства и формирование структурированных
заявок на обслуживание, что создаёт основу для практического внедрения
разработанного решения в деятельность сервисных организаций.

\section{Результаты тестирования системы, анализ корректности
интеллектуальной обработки обращений и особенности
развертывания}\label{cha:chapter3}

\subsection{Результаты тестирования
системы}\label{sec:chapter3-tests}

Помимо автоматизированных проверок выполнена ручная приёмка по требованию
TR-007 (usability, мобильное представление, сценарии ролей клиента, менеджера
и администратора). Протокол зафиксирован в документе
\texttt{docs/manual-acceptance-tr007.md}; по всем пяти сценариям на
настольном и мобильном (375\,px) viewport получен результат Pass.
Критических дефектов, блокирующих демонстрацию и пилотную эксплуатацию,
не выявлено.

Тестирование разработанной системы проводилось на уровне интеграционных
сценариев, охватывающих полный цикл первичной консультации клиента:
инициацию сессии, ввод диагностической информации, обработку обращения
серверной частью, формирование предварительного результата и передачу
заявки менеджеру. Проверки выполнялись в конфигурации, соответствующей
рабочему контуру, что позволило оценить работоспособность решения в
условиях, максимально приближенных к эксплуатации.

Проведенные испытания подтвердили корректность функционирования ролевой
модели доступа, устойчивость обработки пользовательских сообщений и
согласованность данных при сохранении консультационных сессий и заявок.
Отдельно контролировалась корректность сценариев административного
управления контентом и пользователями, а также обработка конкурентных
операций при параллельной работе с заявками. Модульные тесты адаптера
ответа LLM (\texttt{backend/tests/unit/ai-adapter.test.js}) подтверждают
корректность нормализации JSON и устойчивость при недоступности модели
(FR-025b). Нагрузочные испытания выполнялись скриптом
\texttt{tests/perf/k6-consultation.js} с контролем порога p95 времени
ответа. Критических ошибок, препятствующих практической эксплуатации
системы, в ходе проверки не выявлено.

\subsection{Анализ корректности интеллектуальной обработки
обращений}\label{sec:chapter3-ai}

Корректность интеллектуальной обработки анализировалась на типовых
пользовательских запросах с различной степенью полноты исходных данных.
Оценка проводилась по содержательной релевантности формируемых
диагностических гипотез, логической согласованности рекомендаций и
пригодности структурированного результата для дальнейшей работы менеджера
автосервиса.

По результатам анализа установлено, что модуль интеллектуальной обработки
стабильно выделяет ключевые признаки обращения и формирует ответ в
формате, пригодном для включения в регламент первичной консультации.
Применение локальной модели обеспечивает независимость процесса обработки
от внешних сервисов и снижает риски утечки данных. При этом
интеллектуальный вывод используется как предварительная оценка и не
заменяет очную техническую диагностику автомобиля в сервисном центре.

Фрагмент серверной маршрутизации обращения к модулю интеллектуальной
обработки представлен в листинге \ref{lst:consult-route}.

\begin{lstlisting}[caption={Маршрут обработки запроса первичной консультации},label={lst:consult-route}]
router.post('/consultations', async (req, res, next) => {
  try {
    const payload = {
      vehicle: req.body.vehicle,
      symptomText: req.body.symptomText,
      sessionId: req.body.sessionId,
    };
    const result = await consultationsService.process(payload);
    return res.status(200).json(result);
  } catch (error) {
    return next(error);
  }
});
\end{lstlisting}

\subsection{Особенности развертывания
системы}\label{sec:chapter3-deploy}

Развертывание системы организовано в воспроизводимой конфигурации,
включающей серверное приложение, базу данных и локальный интеллектуальный
модуль. Практический контур ввода решения в эксплуатацию включает запуск
инфраструктурных сервисов, применение миграций базы данных, инициализацию
справочных данных, запуск серверного процесса и контроль прохождения
критического пользовательского сценария.

Выбранный подход обеспечивает технологическую однородность среды
разработки и эксплуатации, упрощает сопровождение и позволяет масштабиро
вать систему без изменения базовой архитектуры. Использование централизо
ванного хранилища данных и модульной серверной структуры формирует
устойчивую основу для дальнейшего расширения функциональности.

\subsection{Выводы по третьей
главе}\label{sec:chapter3-summary}

В третьей главе подтверждена практическая работоспособность разработанной
системы в части первичной консультации клиентов автомобильного сервиса.
Результаты интеграционного тестирования показали корректное выполнение
ключевых функциональных сценариев, включая обработку обращений,
формирование предварительного результата и передачу заявки в рабочий
контур менеджера.

Анализ интеллектуальной обработки показал, что подсистема формирует
устойчивый структурированный результат, пригодный для использования в
процессе первичной консультации. Рассмотренные особенности развертывания
подтвердили готовность решения к эксплуатации в прикладном контуре и его
перспективность для дальнейшего развития.

\section*{Список использованных источников}
\addcontentsline{toc}{section}{Список использованных источников}
\begin{thebibliography}{99}
\bibitem{src:iso25010}
ISO/IEC 25010:2011. Systems and software engineering -- Systems and software quality requirements and evaluation (SQuaRE) -- System and software quality models. -- Geneva: ISO, 2011.
\bibitem{src:owasp}
OWASP Foundation. OWASP Top 10: The Ten Most Critical Web Application Security Risks [Электронный ресурс]. -- 2021. -- URL: https://owasp.org/www-project-top-ten/ (дата обращения: 07.05.2026).
\bibitem{src:nist}
NIST AI Risk Management Framework (AI RMF 1.0) [Электронный ресурс]. -- National Institute of Standards and Technology, 2023. -- URL: https://www.nist.gov/itl/ai-risk-management-framework (дата обращения: 07.05.2026).
\bibitem{src:rest}
Fielding R. T. Architectural Styles and the Design of Network-based Software Architectures. -- Irvine: University of California, 2000.
\bibitem{src:docker}
Docker Documentation [Электронный ресурс]. -- URL: https://docs.docker.com/ (дата обращения: 07.05.2026).
\bibitem{src:postgres}
PostgreSQL Documentation [Электронный ресурс]. -- URL: https://www.postgresql.org/docs/ (дата обращения: 07.05.2026).
\bibitem{src:node}
Node.js Documentation [Электронный ресурс]. -- URL: https://nodejs.org/docs/latest/api/ (дата обращения: 07.05.2026).
\end{thebibliography}

\end{document}
